%% kapa: exam=1 round=21 year=2010 period=2 form=A subject=appraisal_law subject_ko=부동산관계법규 seq=26 number=66
%% source: Q-Net 감정평가사 기출문제 (cst00309) artlSeq=5024181 file=extracted/1008/5024181/2010년 제21회 감정평가사 1차시험 기출문제/2교시.hwp
%% conversion: kordoc 4.24.51 -> tools/gen1.py (auto); review=proofread-gpt-sol; flags=none
\begin{question}{66}
국유재산법령상 행정재산의 사용허가에 관한 설명으로 옳은 것은?
\begin{choices}
  \item 행정재산의 사용허가기간이 끝난 재산에 대하여 대통령령이 정하는 경우를 제외하고는 10년을 초과하지 아니하는 범위에서 종전의 사용허가를 갱신할 수 있다.
  \item 관리청은 사용허가한 행정재산을 국가나 지방자치단체가 직접 공용이나 공공용으로 사용하기 위하여 필요하게 된 경우에는 그 허가를 취소할 수 있다.
  \item 행정재산을 직접 공용ㆍ공공용 또는 비영리공익사업용으로 사용하려는 지방자치단체에 사용허가하는 경우에는 사용료를 면제할 수 없다.
  \item 관리청은 동일인이 같은 행정재산을 6개월을 초과하여 계속 사용ㆍ수익하는 경우로서 대통령령으로 정하는 경우에는 사용료를 조정할 수 있다.
  \item 관리청은 행정재산을 사용허가하려는 경우, 사용허가의 목적ㆍ성질ㆍ규모 등을 고려하여 필요하다고 인정되면 대통령령으로 정하는 바에 따라 참가자의 자격을 제한하거나 참가자를 지명하여 경쟁에 부치거나 수의(隨意)의 방법으로 할 수 있다.
\end{choices}
\end{question}
\answer{5}
